Specify absorption, reflection and restarting before writing a walk's equations. If direction has memory or the boundary moves with past records, include direction or drawdown in the state. Hitting probabilities satisfy harmonic equations, whereas mean durations satisfy first-step equations with an added time cost. Conditioning on success reweights transitions by subsequent success probabilities; merely deleting the unsuccessful endpoint does not define the conditioned walk.

Martingales offer another route to boundary calculations, but optional stopping needs hypotheses. First stop at $\tau\wedge n$, then justify passage to the limit through uniform integrability or an integrable dominating variable. Almost-sure termination alone is insufficient. Jumps may also overshoot a target, so use the actual terminal value rather than automatically substituting the boundary level.
\par\smallskip
